% xtufte-book — a Unicode-engine (lualatex/xelatex) fork of the classic
% Tufte-LaTeX book class, extending the already-shipped tufte-style-book
% to OpenType/TrueType fonts. Ships with TeX Live; not bundled.
%
% backend=bibtex forced explicitly: the class defaults to biblatex's
% biber backend, which this catalogue's compile loop does not run
% (classic bibtex only).
%
% Simplified heavily from the class's own doc/sample-book.tex, which
% relies on features this pipeline cannot use: a filecontents block
% generating a helper file, `minted-code` (needs -shell-escape), and a
% remote bibliography fetched over the network at compile time
% (\addbibresource[location=remote]{https://...}) — none of which work
% in a sandboxed, offline compile.
%
% Class documentation: texdoc xtufte | https://ctan.org/pkg/xtufte
% Licence: base tufte-latex is Apache-2.0 (Kevin Godby, Bil Kleb, Bill
% Wood); the xtufte extension is LPPL-1.3, (c) 2025 Herbert Voß.
\PassOptionsToPackage{backend=bibtex}{biblatex}
\documentclass[justified,notoc]{xtufte-book}

\usepackage{booktabs}

\title[Notes on Attention]{Notes on Attention}
\author[Owen Marsh]{Owen Marsh}
\publisher{}

\addbibresource{references.bib}

\begin{document}

\maketitlepage

\tableofcontents

\chapter{What Attention Actually Costs}

Herbert Simon observed in 1971 that ``a wealth of information creates
a poverty of attention''\sidenote{Simon was writing about information
overload decades before the internet made the problem
visceral.}\cite{simon1971designing}. The claim sounds obvious now, but
its consequence is not: if attention is the scarce resource, then
every design decision that competes for it has a real cost, whether or
not that cost shows up on any balance sheet.

\newthought{Notification design} is the clearest modern instance of
this. A single badge on an app icon is a standing claim on your future
attention — it will be resolved eventually, and until it is, some
small fraction of your working memory is reserved for it, whether or
not you are consciously thinking about it.

\begin{marginfigure}
\caption{A schematic: each open notification is a small, ongoing tax
on working memory, not a one-time cost.}
\end{marginfigure}

\section{The Interruption Tax}

Gloria Mark's research on workplace interruption found that it takes
an average of 23 minutes to fully return to a task after an
interruption\sidenote{The commonly cited figure is 23 minutes and 15
seconds, from continuous logging of office workers'
activity.}\cite{mark2008cost}. This is not because the interrupted task
was hard — it is because switching contexts has its own fixed cost,
paid again on the way back.

\begin{table}
\centering
\caption{Rough interruption costs by type, illustrative only.}
\begin{tabular}{@{}lr@{}}
\toprule
Interruption type & Typical recovery time \\
\midrule
Email notification & 1--2 minutes \\
Chat message & 2--5 minutes \\
Phone call & 10--15 minutes \\
Full context switch (meeting) & 20+ minutes \\
\bottomrule
\end{tabular}
\end{table}

\chapter{Designing for Attention}

If attention is scarce, the natural design response is to spend it
deliberately rather than by default. This does not mean removing all
notifications — some genuinely warrant an interruption — but it means
treating every claim on attention as a decision with a cost, not a
free action.

\printbibliography

\end{document}
